% =============================================================================
%  Security analysis.  Slots into the restructured paper as Section 6 (2.5 pp).
%  The long-form hybrid sequence lives in sections/appendix-hybrids.tex.
% =============================================================================
\section{Security Analysis}
\label{sec:security}

We work in the semi-honest model with static corruptions. Throughout, $\lambda$ is the
computational security parameter, $\kappa$ the statistical one, $\mathbb{F}$ a field with
$|\mathbb{F}| \ge 2^{\kappa + \log N}$, and $\mathsf{negl}$ an unspecified negligible
function. Section~\ref{sec:atrest} then steps outside that model, because the most
serious weakness we identify is not a protocol weakness at all.

\subsection{The functionality}
\label{sec:functionality}

Prior work in this line proves security for \emph{one} client set against \emph{one}
labelled server set. NeuralPSI runs $N$ record instances sharing a receiver polynomial
and a PRF key, so the composite object needs to be named before anything can be proved
about it.

\begin{figure}[htbp]
\fbox{\parbox{0.97\linewidth}{
\textbf{Functionality} $\mathcal{F}_{\mathsf{NeuralPSI}}^{d,T,w,\theta}$\\[2pt]
\textbf{Parties.} A client $\mathsf{C}$ and a server $\mathsf{S}$.\\
\textbf{Parameters.} Code length $d$, sub-sample count $T$, sub-sample width $w$,
acceptance threshold $\theta$, and a public mask family
$M = \{m_j\}_{j\in[T]}$, $m_j \in \{0,1\}^d$, $\|m_j\|_1 = w$.\\[3pt]
\textbf{Inputs.} $\mathsf{C}$ inputs $\tilde{\mathbf{x}} \in \{0,1\}^d$.
$\mathsf{S}$ inputs a database $D = \{(\tilde{\mathbf{y}}_i, \ell_i)\}_{i\in[N]}$ with
$\tilde{\mathbf{y}}_i \in \{0,1\}^d$ and labels $\ell_i \in \mathbb{F}$.\\[3pt]
\textbf{Computation.} For each $i \in [N]$ let
\[
S_i \;=\; \bigl\{\, j \in [T] \;:\; \tilde{\mathbf{x}} \wedge m_j
   \;=\; \tilde{\mathbf{y}}_i \wedge m_j \,\bigr\},
\qquad
\mathsf{acc}(i) \;=\; \bigl[\,|S_i| \ge \theta\,\bigr].
\]
\textbf{Outputs.} $\mathsf{C}$ receives $\mathsf{out} =
\{\,(i, \ell_i) : \mathsf{acc}(i)\,\}$; $\mathsf{S}$ receives $\bot$.
}}
\caption{The composite functionality. Note that the output is a \emph{set}: the scan runs
over all $N$ records and more than one may satisfy the predicate. This differs from the
single-match phrasing used in earlier drafts of this work, and the difference is
observable --- $|\mathsf{out}| > 1$ occurs at a rate we measure in
Section~\ref{sec:params}.}
\label{fig:functionality}
\end{figure}

Two features of Figure~\ref{fig:functionality} deserve emphasis. First, acceptance is a
deterministic function of $\tilde{\mathbf{x}}$, $\tilde{\mathbf{y}}_i$ and the public mask
family, so it depends on the pair only through which positions differ; because the masks
are drawn uniformly it depends in distribution only on the Hamming weight
$H = \|\tilde{\mathbf{x}} \oplus \tilde{\mathbf{y}}_i\|_1$. Every error rate in
Section~\ref{sec:params} follows from that observation. Second, the functionality says
nothing about \emph{authentication}: no party other than $\mathsf{C}$ learns anything, so
a relying party cannot act on the result. We return to this in
Section~\ref{sec:whoauth}.

\paragraph{The masks are public.} We fix this now because the leakage analysis depends on
it. Taking $M$ public is the conservative reading: it gives the adversary strictly more
than a server-private mask family would, and our claims hold under it. Making $M$
server-private would shrink $\mathcal{L}_{\mathsf{C}}$ below, but the protocol as
constructed delivers the collision pattern to the client, so a private $M$ would require a
structural change rather than a parameter choice.

\subsection{The leakage function}
\label{sec:leakage}

$\mathcal{F}_{\mathsf{NeuralPSI}}$ is what we would like to realise. What the protocol
realises is $\mathcal{F}_{\mathsf{NeuralPSI}}$ \emph{augmented with leakage}:
\begin{align}
\mathcal{L}_{\mathsf{C}}(\tilde{\mathbf{x}}, D) &=
  \Bigl(\;N,\;\; \bigl\{\,(i, S_i) \;:\; |S_i| \ge \theta \,\bigr\} \;\Bigr),
  \label{eq:leakC}\\[2pt]
\mathcal{L}_{\mathsf{S}}(\tilde{\mathbf{x}}, D) &= \bot .
  \label{eq:leakS}
\end{align}

Equation~\eqref{eq:leakC} is strictly more than ``the identity of the matched record''.
The client learns \emph{which} sub-samples collided, and since $M$ is public, each
$j \in S_i$ pins $w$ known positions of $\tilde{\mathbf{y}}_i$ to known values. The client
also learns $N$, the number of records tested. We quantify the consequences in
Section~\ref{sec:attacks}; the short version is that $\mathcal{L}_{\mathsf{C}}$ is enough
to reconstruct a matched record outright.

\subsection{Security theorem}

\begin{theorem}[Semi-honest security]
\label{thm:main}
Let $F$ be a PRF, and instantiate the sub-protocols with a masked OPRF
$\Pi_{\mathsf{mOPRF}}$ and a vector ring-OLE $\Pi_{\mathsf{VROLE}}$, each semi-honest
secure. Then the NeuralPSI protocol of Section~\ref{sec:design} securely realises
$\mathcal{F}_{\mathsf{NeuralPSI}}^{d,T,w,\theta}$ with leakage
$(\mathcal{L}_{\mathsf{C}}, \mathcal{L}_{\mathsf{S}})$ against a semi-honest adversary
statically corrupting one party. Concretely, for every such adversary $\mathcal{A}$ there
is a simulator whose output is distinguishable from $\mathcal{A}$'s view with advantage at
most
\begin{equation}
\label{eq:bound}
\varepsilon \;\le\;
   \varepsilon_{\mathsf{mOPRF}}
 + \varepsilon_{\mathsf{VROLE}}
 + \varepsilon_{\mathsf{LPN}}
 + \underbrace{\varepsilon_{F}\bigl(q + T\!\cdot\!2^{w}\bigr)}_{\text{PRF, extended domain}}
 + \underbrace{\frac{N(2T+1)}{|\mathbb{F}|}}_{\text{statistical}} ,
\end{equation}
where $q$ is the adversary's own query count.
\end{theorem}

Two terms in~\eqref{eq:bound} are easy to omit and both matter. The PRF term is charged
for $T \cdot 2^{w}$ extra oracle queries, not $q$ alone: each mask restricts the PRF input
to a domain of size $2^{w} = 2^{14}$, so the full input table has only
$T \cdot 2^{w} = 2^{20}$ entries and the reduction must budget for enumerating it. The
statistical term is the price of the polynomial masking step (hybrid $\mathsf{H}_4$
below) and is the reason $|\mathbb{F}|$ must be chosen with $N$ in view; we report the
concrete field size in Section~\ref{sec:eval}.

Both simulators are constructed in Appendix~\ref{app:hybrids}. The corrupt-client simulator
takes $\xt$ and $(\mathsf{out}, \mathcal{L}_{\mathsf{C}})$ and produces the mOPRF outputs, the
sender-side silent-VOLE transcript and the polynomials $\{\hat z_i\}$, sampling uniformly at every
point not pinned by $\mathcal{L}_{\mathsf{C}}$. The corrupt-server simulator has only the
transcript to produce, since the server outputs $\bot$; its one input-dependent message is
$\Delta U = \Delta u(X)$. Worth stating precisely, because the natural argument is slightly wrong:
$\Delta U$ is \emph{not} uniform on $\mathbb{F}^{2T+1}$ but lies in the $(T{+}1)$-dimensional image
of the degree-$\le T$ evaluation map, so the simulator samples a uniform polynomial of degree
$\le T$ and outputs its evaluation vector. Since $u'_0$ is uniform, $u_0 - u'_0$ is uniform
\emph{in that subspace} and the indistinguishability is statistical, not computational.

\subsection{Hybrids}
\label{sec:hybrids}

$\mathsf{H}_0$ is the real execution. $\mathsf{H}_1$ replaces $\Pi_{\mathsf{mOPRF}}$ by
$\mathcal{F}_{\mathsf{mOPRF}}^{T}$ (semi-honest mOPRF security); $\mathsf{H}_2$ replaces $F_k$ by
a random function (PRF security, charged at $q + T\cdot2^{w}$ queries as above); $\mathsf{H}_3$
replaces $\Pi_{\mathsf{VROLE}}$ by $\mathcal{F}_{\mathsf{VROLE}}$ (base-VOLE security and LPN).

$\mathsf{H}_4$ is the load-bearing \emph{statistical} step, and is the one most easily left
unstated: for every non-colliding point $x_j$ it replaces
$v_1(x_j) = f(x_j) + \bigl(\prod_i (x_j - \yt_i)\bigr)\cdot R(x_j)$ by a uniform element of
$\mathbb{F}$. This holds because $R$ is uniform of degree $\le T$ and the evaluation map on any
set of at most $T+1$ points is a bijection onto $\mathbb{F}^{T+1}$, so the masking terms are
jointly uniform --- \emph{provided}
\begin{equation}
\label{eq:h4}
\bigl|\{\, j \in [T] : j \notin S_i \,\}\bigr| \;\le\; \deg R + 1 \;=\; T+1,
\end{equation}
which holds unconditionally here since $|[T]| = T$, but is a design constraint rather than an
accident. Over $N$ records and $2T+1$ points the union bound gives the $N(2T+1)/|\mathbb{F}|$
term in~\eqref{eq:bound}. Finally $\mathsf{H}_5$ replaces $\ell'_i = F_k(\ell_i)$ by uniform
elements (immediate from $\mathsf{H}_2$), and $\mathsf{H}_6 = \mathsf{Sim}_{\mathsf{C}}$ is
identical by inspection.

\paragraph{Why $N$ instances do not compose for free.}
The step that a citation cannot discharge is the reuse across records. All $N$ record
instances share one receiver polynomial $u_0$ --- that sharing is precisely what makes the
vector ring-OLE batching asymmetric and cheap --- and one PRF key $k$, across all records
\emph{and} all sessions. Hybrids $\mathsf{H}_2$ and $\mathsf{H}_4$ are therefore stated
over the whole database rather than per record: $\mathsf{H}_2$ replaces the single shared
$F_k$ once, which is why its cost appears once rather than $N$ times, and $\mathsf{H}_4$
is applied simultaneously at all $N(2T+1)$ points with the union bound
in~\eqref{eq:bound}. Constraint~\eqref{eq:h4} must hold for every record at once; it does,
but only because $\deg R \ge T$, which is a design requirement rather than an accident and
is easy to violate when tuning $T$ downward for communication.

The long-form hybrids, with the distinguisher reductions written out, are in
Appendix~\ref{app:hybrids}.

\subsection{The stored template is a bearer credential}
\label{sec:atrest}

Theorem~\ref{thm:main} covers the protocol and says nothing about data at rest, which is where
this design is weakest --- strictly weaker than the homomorphic baseline, whose stored feature
vectors are encrypted under a key the server never holds.

The server stores $(\yt_i, \ell_i)$ in the clear. An adversary who reads $\yt_i$ and submits it
as a query obtains $H = 0$, hence $|S_i| = T$ and acceptance with probability $1$. \emph{No
inversion is required}: the stored code \emph{is} a working credential, so a database breach is
immediate and total impersonation of every enrolled user.

Inversion is an additional problem. The quantiser is over-determined --- 128 hyperplanes cut
$\mathbb{S}^{15}$ into $2\sum_{k<16}\binom{127}{k} \approx 2^{64.5}$ cells --- so it is not
one-way in any useful sense. Its intrinsic angular resolution is
$\text{cells}^{-1/15} \approx 2.9^\circ$, and the standard one-bit compressed-sensing heuristic
$k/m = 16/128$ gives $7.2^\circ$, against a measured genuine intra-class spread of
$\pi\cdot 8.3/128 = 11.7^\circ$. Reconstruction is thus $1.6$--$4\times$ \emph{tighter} than
genuine variation: the recovered direction re-quantises inside the accept region and transfers
to any system using the same extractor.

\begin{table}[htbp]
\centering
\caption{The stored template against ISO/IEC 24745~\cite{ISO24745}. NeuralPSI as
specified fails all three requirements.}
\label{tab:iso24745}
\small
\begin{tabular}{@{}p{0.19\linewidth}p{0.10\linewidth}p{0.62\linewidth}@{}}
\toprule
\textbf{Requirement} & \textbf{Status} & \textbf{Why} \\
\midrule
Irreversibility & fails & Twice: direct replay of the stored code, and one-bit
compressed-sensing recovery of the embedding direction to $2.9$--$7.2^\circ$. \\
Unlinkability & fails & $(P, \mu, \tau)$ and the projection seed are public and fixed, so
one finger yields the same 128-bit code in every deployment and codes are directly
Hamming-comparable across databases. There is no diversification parameter. \\
Renewability & fails & A second independent code cannot be issued from the same finger.
This is the requirement the application most needs, since a compromised fingerprint
cannot be reissued. \\
\bottomrule
\end{tabular}
\end{table}

\paragraph{Mitigations, and what they actually buy.}
\emph{(i) A client-side per-enrolment XOR mask}: the device holds $r_u \in \{0,1\}^{d}$ and the
server stores $\yt \oplus r_u$. XOR preserves Hamming distance exactly, so \textbf{accuracy is
unchanged} --- every number in \S\ref{sec:eval} carries over --- and 16 bytes of client state
buys unlinkability and renewability. It does \emph{not} stop replay of a stolen record.
\emph{(ii) Envelope encryption at rest}, key in an HSM, decrypting to memory for the scan:
unglamorous, orthogonal to the protocol, sound in the semi-honest model, and what we would
deploy. \emph{(iii)} We explicitly reject one tempting option --- storing the mOPRF images
$\{F_k(\yt_i \wedge m_j)\}_j$ instead of the code. It is keyed and locality-preserving and the
server computes those values anyway, but the per-mask domain is only $2^{w} = 2^{14}$, so anyone
who can evaluate $F_k$ inverts each image by exhaustive search; repairing it needs a per-record
salt that destroys the equality-preservation the protocol depends on.

Applying (i) and (ii) yields unlinkability and renewability. \textbf{Full at-rest
confidentiality against a compromised server remains out of scope}, and this posture is weaker
than the encrypted-template baseline's --- a limitation we state rather than resolve.

\subsection{Who is actually authenticated}
\label{sec:whoauth}

As specified, no relying party learns the outcome, so no session can be granted:
$\mathcal{F}_{\mathsf{NeuralPSI}}$ is a private identification lookup, not an authentication
protocol. The credential step of Section~\ref{sec:design} closes the gap for one MAC, one nonce
and 32 bytes, with the nonce and session id making $\sigma$ non-replayable --- without them one
intercepted transcript is a permanent credential, since $\ell$ is long-lived.

That step has a consequence for how Section~\ref{sec:params} should be read. A client cannot
forge $\sigma$ without $\ell_i$, so system security reduces to recovering \emph{some} $\ell_i$,
which is exactly the false-accept question. In a conventional matcher a false accept is one bad
login; here it hands the attacker a \textbf{permanent per-user credential}. \emph{The
false-accept rate is a credential-compromise rate}, which is why we treat it as a security
parameter rather than an accuracy metric.

\subsection{Attacks}
\label{sec:attacks}

\paragraph{(a) Impersonation without search.}
The parameters were chosen against an attacker submitting \emph{uniformly random} codes,
for which the per-record accept probability is $q_{\mathsf{rand}} = 2.954\times10^{-5}$.
But the cheapest attack is not to guess: it is to present one's own finger. Measured over
$2.88\times10^{6}$ held-out impostor pairs, the per-record accept probability for a
\emph{real but wrong} finger is $q_{\mathsf{bio}} = 4.578\times10^{-2}$ --- $1{,}550\times$
larger. Expected queries to impersonate somebody, $1/(1-(1-q)^N)$:

\begin{center}\small
\begin{tabular}{@{}lrrrr@{}}
\toprule
$N$ & $512$ & $5{,}000$ & $10^{5}$ & $10^{6}$ \\
\midrule
random code ($q_{\mathsf{rand}}$) & $67$ & $7$ & $1.05$ & $1$ \\
\textbf{zero-effort finger} ($q_{\mathsf{bio}}$) & $\mathbf{1}$ & $\mathbf{1}$ &
  $\mathbf{1}$ & $\mathbf{1}$ \\
\bottomrule
\end{tabular}
\end{center}

At every database size we measured, a single presentation of an unrelated finger
impersonates \emph{someone} with probability above $99\%$. Rate limiting does not
mitigate an attack that succeeds on the first attempt.

\paragraph{(b) Template recovery from the collision pattern.}
On any accept, $\mathcal{L}_{\mathsf{C}}$ gives the client $S_i$. With $\theta = 2$ the
client can enumerate all $\binom{64}{2} = 2{,}016$ candidate pairs, interpolate, and test,
so it learns $S_i$ exactly and cheaply. Because the masks are public, each $j \in S_i$
pins $w = 14$ known positions of $\tilde{\mathbf{y}}_i$. At $H \approx 8$ a genuine query
produces $T \cdot q(8) = 64 \times 0.38493 \approx 24.6$ collisions, covering
$24.6 \times 14 \approx 345$ position-slots over $128$ positions: \textbf{one successful
authentication determines the target record}. Recovering a \emph{chosen} user costs
$1/q_{\mathsf{bio}} \approx 22$ queries at the measured rate; recovering \emph{a} user's
full template at $N \ge 10^{5}$ costs one.

\paragraph{(c) The mOPRF dictionary.}
Each mask restricts the PRF domain to $2^{w} = 2^{14}$, and the entire table is
$T \cdot 2^{w} = 2^{20}$. A client learns one entry per mask per session, so a
coupon-collector argument gives roughly $1.1\times10^{4}$ sessions for $50\%$ coverage per
mask --- already sufficient, since 32 of 64 masks at 14 bits each over-determines a
128-bit code --- and about $1.7\times10^{5}$ sessions for full coverage. ``PRF-masked''
thus provides at most $w = 14$ bits of preimage resistance per sub-sample against a
querying client.

\paragraph{Required deployment controls.} Consequently: (i) rate limiting per account
\emph{and} per relying party, with the explicit admission that it does not mitigate (a) at
large $N$; (ii) authenticated enrolment, so an attacker cannot enrol chosen codes to probe
the space; (iii) per-epoch rotation of $k$ and of the mask family, at the stated
re-enrolment cost; (iv) sensor attestation as an explicit deployment assumption; and
(v) the nonce binding of Section~\ref{sec:whoauth}.

\paragraph{What we do not claim.} Malicious security --- cut-and-choose techniques exist for
related structure-aware PSI settings~\cite{GarimellaEtAl23}, but they target geometric-ball
rather than threshold predicates, so their overheads do not transfer and we do not estimate
them. At-rest confidentiality against a compromised server. That the quantiser is one-way: it
demonstrably is not, and no part of Theorem~\ref{thm:main} relies on it. And that the measured
false-accept rates are acceptable at the sizes we tested --- \S\ref{sec:params} shows they are
not, and why.
